

\documentclass[12pt]{article}
\usepackage{amsmath}
\usepackage{hyperref}
\date{}
\begin{document}
\begin{titlepage}
  \title{Zero Coupon Yield Curve for India: methodology and sources}
  \author{Finance Research Group}
  \maketitle
\end{titlepage}
\section*{Panel 1: Daily ZCYC (Jun 2001 - Sep 2020) }
Panel 1 is a video of the daily Zero Coupon Yield Curve (ZCYC) in India for the period Jun 2001 - Sep 2020. 
The ZCYC parameters are modeled using the \href{https://econpapers.repec.org/article/ucpjnlbus/v_3a60_3ay_3a1987_3ai_3a4_3ap_3a473-89.htm}{Nelson-Siegel model (3-factor model)} for the period Jun 2001 - Dec 2010 and using the \href{https://econpapers.repec.org/paper/nbrnberwo/4871.htm}{Nelson-Siegel-Svensson (5-factor model)} for the period Jan 2011 - Sep 2020. 
The Nelson-Siegel model parameters for the Jun 2001 - Dec 2010 period are from NSE and that for the Jan 2011 - Sep 2020 period, are from CCIL. \\ \\

We calculate the ZCYC for daily maturities from 1 day to 30 years using the following formulae - 
\begin{itemize}
\item For the period Jun 2001 - Dec 2010, we use Nelson-Siegel (3-factor model)
  \begin{equation}
    R(m)= \beta_0 + (\beta_1 + \beta_2)\cdot\frac{[1-\exp(-\frac{m}{\tau})]}{\frac{m}{\tau}} - \beta_2\cdot[\exp(-\frac{m}{\tau})] 
  \end{equation}
  where, m is maturity,\\
  $\beta_0$ , $\beta_1$, $\beta_2$ and $\tau$ are parameters of the equation,\\
  and, R(m) is the yield for maturity (m).
  
\item For the period Jan 2011 - Sep 2020, we use Nelson-Siegel-Svensson (5-factor model) 
  
  \begin{equation}
\begin{split}
  R(m)= \beta_0 + (\beta_1 + \beta_2)\cdot[\frac{[1-\exp(-\frac{m}{\tau_1})]}{\frac{m}{\tau_1}}]  - \beta_2\cdot[\exp(-\frac{m}{\tau_1})]\\
  + \beta_3\cdot[\frac{[1-\exp(-\frac{m}{\tau_2})]}{\frac{m}{\tau_2}}]  - \beta_3\cdot[\exp(-\frac{m}{\tau_2})] 
\end{split}
\end{equation}
  where, m is maturity,\\
  $\beta_0$ , $\beta_1$, $\beta_2$, $\beta_3$, $\tau_1$ and $\tau_2$ are parameters of the equation,\\
  and, R(m) is the yield for maturity (m).

\end{itemize} 

The daily ZCYC is colour coded to capture the inversion of the term structure using the 10 year and the 91 day maturity.
\begin{itemize}
\item When 10 year maturity is $>=$ 91 day maturity, the ZCYC curve is blue.
\item When 10 year maturity is $<$ 91 day maturity, the ZCYC curve is red.
\end{itemize}

\section*{Panel 2: Evolution of ZCYC (Jun 2001 - Sep 2020)}

It compares the ZCYC over time by showing one ZCYC for each year, which is as at the last day of March in that particular year. 
For example - for the year 2011, the curve represents ZCYC on $31^{st}$ March 2011. The computation of ZCYC is as per the method described for Panel 1.

\section*{Panel 3: Term spread}

The term spread is calculated for 2 maturity pair :
\begin{itemize}
\item 10 year and 91 day, and 
\item 10 year and 2 year. 
\end{itemize}
The term spread is computed at a daily frequency as the rolling window average of the last 30 trading day spread.
\end{document}
